% preamble_common_en.tex
% English counterpart of preamble_common.tex. Shared preamble for all print
% editions of the English book (main_en.tex, main_kdp_print_en.tex,
% main_bod_print_en.tex). Holds everything that does NOT depend on the trim
% size. \documentclass and \geometry stay in each wrapper file, since they
% differ per edition.

\usepackage{iftex}
\usepackage[english]{babel}
\usepackage[T1]{fontenc}
\usepackage{textcomp}
\ifPDFTeX
    \usepackage[utf8]{inputenc}
\fi
\usepackage{lmodern}

% protrusion=false: by default, microtype lets single characters (quotation
% marks, periods, hyphens, often also italic glyphs such as \emph) protrude
% a tiny amount (a fraction of a point) PAST the actual text block margin
% ("hanging punctuation", an intentional, common typographic feature). KDP's
% automatic check, however, appears to flag any detected content beyond the
% text box as "outside the margins", regardless of how much safety margin
% actually remains to the trim edge. KDP itself documents exactly this
% failure case for italic text (source: kdpcommunity.com, thread "Continued
% formatting error/text outside margins", as of 2026-07). protrusion only
% affects fractions of a point and barely changes line breaking - expansion
% (letter-spacing adjustment) stays active and still delivers most of
% microtype's benefit for justified text.
\usepackage[protrusion=false]{microtype}
\usepackage{geometry}

% Extra justification slack: catches cases that microtype alone can no
% longer balance (long English compound terms, narrow columns on smaller
% trims). Without these values, LaTeX produces a visibly overfull hbox
% that sticks out past the text block in such cases, instead of breaking
% the line cleanly with slightly more word spacing.
\tolerance=1000
\emergencystretch=3em

\usepackage{graphicx}
\usepackage{adjustbox}
\usepackage{booktabs}
\usepackage{longtable}
\usepackage{array}
\usepackage{xcolor}

% Print editions (KDP, BoD) set \druckausgabetrue BEFORE % preamble_common_en.tex
% English counterpart of preamble_common.tex. Shared preamble for all print
% editions of the English book (main_en.tex, main_kdp_print_en.tex,
% main_bod_print_en.tex). Holds everything that does NOT depend on the trim
% size. \documentclass and \geometry stay in each wrapper file, since they
% differ per edition.

\usepackage{iftex}
\usepackage[english]{babel}
\usepackage[T1]{fontenc}
\usepackage{textcomp}
\ifPDFTeX
    \usepackage[utf8]{inputenc}
\fi
\usepackage{lmodern}

% protrusion=false: by default, microtype lets single characters (quotation
% marks, periods, hyphens, often also italic glyphs such as \emph) protrude
% a tiny amount (a fraction of a point) PAST the actual text block margin
% ("hanging punctuation", an intentional, common typographic feature). KDP's
% automatic check, however, appears to flag any detected content beyond the
% text box as "outside the margins", regardless of how much safety margin
% actually remains to the trim edge. KDP itself documents exactly this
% failure case for italic text (source: kdpcommunity.com, thread "Continued
% formatting error/text outside margins", as of 2026-07). protrusion only
% affects fractions of a point and barely changes line breaking - expansion
% (letter-spacing adjustment) stays active and still delivers most of
% microtype's benefit for justified text.
\usepackage[protrusion=false]{microtype}
\usepackage{geometry}

% Extra justification slack: catches cases that microtype alone can no
% longer balance (long English compound terms, narrow columns on smaller
% trims). Without these values, LaTeX produces a visibly overfull hbox
% that sticks out past the text block in such cases, instead of breaking
% the line cleanly with slightly more word spacing.
\tolerance=1000
\emergencystretch=3em

\usepackage{graphicx}
\usepackage{adjustbox}
\usepackage{booktabs}
\usepackage{longtable}
\usepackage{array}
\usepackage{xcolor}

% Print editions (KDP, BoD) set \druckausgabetrue BEFORE % preamble_common_en.tex
% English counterpart of preamble_common.tex. Shared preamble for all print
% editions of the English book (main_en.tex, main_kdp_print_en.tex,
% main_bod_print_en.tex). Holds everything that does NOT depend on the trim
% size. \documentclass and \geometry stay in each wrapper file, since they
% differ per edition.

\usepackage{iftex}
\usepackage[english]{babel}
\usepackage[T1]{fontenc}
\usepackage{textcomp}
\ifPDFTeX
    \usepackage[utf8]{inputenc}
\fi
\usepackage{lmodern}

% protrusion=false: by default, microtype lets single characters (quotation
% marks, periods, hyphens, often also italic glyphs such as \emph) protrude
% a tiny amount (a fraction of a point) PAST the actual text block margin
% ("hanging punctuation", an intentional, common typographic feature). KDP's
% automatic check, however, appears to flag any detected content beyond the
% text box as "outside the margins", regardless of how much safety margin
% actually remains to the trim edge. KDP itself documents exactly this
% failure case for italic text (source: kdpcommunity.com, thread "Continued
% formatting error/text outside margins", as of 2026-07). protrusion only
% affects fractions of a point and barely changes line breaking - expansion
% (letter-spacing adjustment) stays active and still delivers most of
% microtype's benefit for justified text.
\usepackage[protrusion=false]{microtype}
\usepackage{geometry}

% Extra justification slack: catches cases that microtype alone can no
% longer balance (long English compound terms, narrow columns on smaller
% trims). Without these values, LaTeX produces a visibly overfull hbox
% that sticks out past the text block in such cases, instead of breaking
% the line cleanly with slightly more word spacing.
\tolerance=1000
\emergencystretch=3em

\usepackage{graphicx}
\usepackage{adjustbox}
\usepackage{booktabs}
\usepackage{longtable}
\usepackage{array}
\usepackage{xcolor}

% Print editions (KDP, BoD) set \druckausgabetrue BEFORE % preamble_common_en.tex
% English counterpart of preamble_common.tex. Shared preamble for all print
% editions of the English book (main_en.tex, main_kdp_print_en.tex,
% main_bod_print_en.tex). Holds everything that does NOT depend on the trim
% size. \documentclass and \geometry stay in each wrapper file, since they
% differ per edition.

\usepackage{iftex}
\usepackage[english]{babel}
\usepackage[T1]{fontenc}
\usepackage{textcomp}
\ifPDFTeX
    \usepackage[utf8]{inputenc}
\fi
\usepackage{lmodern}

% protrusion=false: by default, microtype lets single characters (quotation
% marks, periods, hyphens, often also italic glyphs such as \emph) protrude
% a tiny amount (a fraction of a point) PAST the actual text block margin
% ("hanging punctuation", an intentional, common typographic feature). KDP's
% automatic check, however, appears to flag any detected content beyond the
% text box as "outside the margins", regardless of how much safety margin
% actually remains to the trim edge. KDP itself documents exactly this
% failure case for italic text (source: kdpcommunity.com, thread "Continued
% formatting error/text outside margins", as of 2026-07). protrusion only
% affects fractions of a point and barely changes line breaking - expansion
% (letter-spacing adjustment) stays active and still delivers most of
% microtype's benefit for justified text.
\usepackage[protrusion=false]{microtype}
\usepackage{geometry}

% Extra justification slack: catches cases that microtype alone can no
% longer balance (long English compound terms, narrow columns on smaller
% trims). Without these values, LaTeX produces a visibly overfull hbox
% that sticks out past the text block in such cases, instead of breaking
% the line cleanly with slightly more word spacing.
\tolerance=1000
\emergencystretch=3em

\usepackage{graphicx}
\usepackage{adjustbox}
\usepackage{booktabs}
\usepackage{longtable}
\usepackage{array}
\usepackage{xcolor}

% Print editions (KDP, BoD) set \druckausgabetrue BEFORE \input{preamble_common_en}.
% KDP's check appears to flag pages with clickable hyperlink annotations
% (cross-references, glossary terms, the whole linked table of contents) as
% suspicious regardless of their color - black link color alone did not
% remove the warnings. For print editions, hyperref is therefore loaded in
% draft mode: all \href/\ref/\gls commands remain usable and still set their
% text normally, but no longer create a PDF link annotation. The working
% edition (main_en.tex) keeps full, clickable links for on-screen use.
\makeatletter
\@ifundefined{ifdruckausgabe}{\newif\ifdruckausgabe\druckausgabefalse}{}
\makeatother
\ifdruckausgabe
    \usepackage[draft]{hyperref}
\else
    \usepackage{hyperref}
\fi

\usepackage{listings}
\usepackage{etoolbox}
\usepackage{needspace}
\usepackage{enumitem}
\usepackage{csquotes}
\usepackage{caption}
\usepackage{float}
\usepackage{tikz}
\usetikzlibrary{positioning,arrows.meta,shadows,calc,quotes,backgrounds}
\usepackage{tcolorbox}
\usepackage{amsmath}
\usepackage{makeidx}
\usepackage[acronym,toc]{glossaries}

\makeindex
\makeglossaries

\hypersetup{
    colorlinks=true,
    linkcolor=blue!50!black,
    urlcolor=blue!50!black,
    citecolor=blue!50!black,
    pdftitle={SysML v2, MBSE, and ROS2 for Beginners},
    pdfauthor={Wolfgang Becker}
}

\RedeclareSectionCommand[tocnumwidth=3.8em]{section}
\RedeclareSectionCommand[tocnumwidth=4.8em]{subsection}
\RedeclareSectionCommand[tocnumwidth=5.8em]{subsubsection}

% Section headings must not stay alone at the bottom of a page. The
% reserved space must fit the heading plus a bit of following text, or the
% start of a short code block.
\pretocmd{\section}{\Needspace{12\baselineskip}}{}{}
\pretocmd{\subsection}{\Needspace{8\baselineskip}}{}{}
\pretocmd{\subsubsection}{\Needspace{7\baselineskip}}{}{}

% Table of contents only down to section level. Subsection entries alone
% (3-digit numbers such as 1.10.1 plus a deeply indented title, not yet
% subsubsection - counting depth: chapter=0, section=1, subsection=2) already
% exceed the TOC line width on the narrower 6x9in trim and run past the
% right margin. tocdepth=2 (the previous setting) still allowed subsections,
% which was not enough - hence tocdepth=1 now. Numbering in the body text is
% unchanged (only \tocdepth, not \secnumdepth); subsections are still
% numbered, they just no longer appear in the table of contents.
\setcounter{tocdepth}{1}

\lstdefinestyle{code}{
    basicstyle=\ttfamily\small,
    breaklines=true,
    frame=single,
    numbers=left,
    numberstyle=\tiny,
    xleftmargin=0.5cm,
    columns=fullflexible,
    showstringspaces=false,
    tabsize=2,
    extendedchars=true,
    inputencoding=utf8,
    upquote=true,
    literate=
        {ä}{{\"a}}1 {ö}{{\"o}}1 {ü}{{\"u}}1
        {Ä}{{\"A}}1 {Ö}{{\"O}}1 {Ü}{{\"U}}1
        {ß}{{\ss{}}}1
}
\lstset{style=code}

% Code blocks must not start with only a few lines at the bottom of a page.
% With less than about eight code lines of space left, the whole block
% moves to the next page; longer listings may still break.
\BeforeBeginEnvironment{lstlisting}{\Needspace{10\baselineskip}}

% Style for plain text/tree displays (directory trees, flow sketches) that
% should look like \verbatim but - unlike \verbatim - must wrap long lines
% so they do not run past the text block on narrow trims (e.g. 6x9in).
\lstdefinestyle{plaintext}{
    style=code,
    numbers=none,
    frame=single,
}

\newtcolorbox{lernziele}{
    colback=blue!3!white,
    colframe=blue!50!black,
    title=Learning Objectives
}
\newtcolorbox{praxisbox}{
    colback=green!3!white,
    colframe=green!40!black,
    title=Practical Example
}
\newtcolorbox{hinweisbox}{
    colback=yellow!8!white,
    colframe=yellow!50!black,
    title=Note
}
\newtcolorbox{uebungbox}{
    colback=gray!5!white,
    colframe=gray!60!black,
    title=Exercise
}

\input{glossary_en}

\title{SysML v2, MBSE, and ROS2 for Beginners\\
\large A Hands-On Textbook Based on an Autonomous Indoor Robot}
\author{Wolfgang Becker}
\date{\today}
.
% KDP's check appears to flag pages with clickable hyperlink annotations
% (cross-references, glossary terms, the whole linked table of contents) as
% suspicious regardless of their color - black link color alone did not
% remove the warnings. For print editions, hyperref is therefore loaded in
% draft mode: all \href/\ref/\gls commands remain usable and still set their
% text normally, but no longer create a PDF link annotation. The working
% edition (main_en.tex) keeps full, clickable links for on-screen use.
\makeatletter
\@ifundefined{ifdruckausgabe}{\newif\ifdruckausgabe\druckausgabefalse}{}
\makeatother
\ifdruckausgabe
    \usepackage[draft]{hyperref}
\else
    \usepackage{hyperref}
\fi

\usepackage{listings}
\usepackage{etoolbox}
\usepackage{needspace}
\usepackage{enumitem}
\usepackage{csquotes}
\usepackage{caption}
\usepackage{float}
\usepackage{tikz}
\usetikzlibrary{positioning,arrows.meta,shadows,calc,quotes,backgrounds}
\usepackage{tcolorbox}
\usepackage{amsmath}
\usepackage{makeidx}
\usepackage[acronym,toc]{glossaries}

\makeindex
\makeglossaries

\hypersetup{
    colorlinks=true,
    linkcolor=blue!50!black,
    urlcolor=blue!50!black,
    citecolor=blue!50!black,
    pdftitle={SysML v2, MBSE, and ROS2 for Beginners},
    pdfauthor={Wolfgang Becker}
}

\RedeclareSectionCommand[tocnumwidth=3.8em]{section}
\RedeclareSectionCommand[tocnumwidth=4.8em]{subsection}
\RedeclareSectionCommand[tocnumwidth=5.8em]{subsubsection}

% Section headings must not stay alone at the bottom of a page. The
% reserved space must fit the heading plus a bit of following text, or the
% start of a short code block.
\pretocmd{\section}{\Needspace{12\baselineskip}}{}{}
\pretocmd{\subsection}{\Needspace{8\baselineskip}}{}{}
\pretocmd{\subsubsection}{\Needspace{7\baselineskip}}{}{}

% Table of contents only down to section level. Subsection entries alone
% (3-digit numbers such as 1.10.1 plus a deeply indented title, not yet
% subsubsection - counting depth: chapter=0, section=1, subsection=2) already
% exceed the TOC line width on the narrower 6x9in trim and run past the
% right margin. tocdepth=2 (the previous setting) still allowed subsections,
% which was not enough - hence tocdepth=1 now. Numbering in the body text is
% unchanged (only \tocdepth, not \secnumdepth); subsections are still
% numbered, they just no longer appear in the table of contents.
\setcounter{tocdepth}{1}

\lstdefinestyle{code}{
    basicstyle=\ttfamily\small,
    breaklines=true,
    frame=single,
    numbers=left,
    numberstyle=\tiny,
    xleftmargin=0.5cm,
    columns=fullflexible,
    showstringspaces=false,
    tabsize=2,
    extendedchars=true,
    inputencoding=utf8,
    upquote=true,
    literate=
        {ä}{{\"a}}1 {ö}{{\"o}}1 {ü}{{\"u}}1
        {Ä}{{\"A}}1 {Ö}{{\"O}}1 {Ü}{{\"U}}1
        {ß}{{\ss{}}}1
}
\lstset{style=code}

% Code blocks must not start with only a few lines at the bottom of a page.
% With less than about eight code lines of space left, the whole block
% moves to the next page; longer listings may still break.
\BeforeBeginEnvironment{lstlisting}{\Needspace{10\baselineskip}}

% Style for plain text/tree displays (directory trees, flow sketches) that
% should look like \verbatim but - unlike \verbatim - must wrap long lines
% so they do not run past the text block on narrow trims (e.g. 6x9in).
\lstdefinestyle{plaintext}{
    style=code,
    numbers=none,
    frame=single,
}

\newtcolorbox{lernziele}{
    colback=blue!3!white,
    colframe=blue!50!black,
    title=Learning Objectives
}
\newtcolorbox{praxisbox}{
    colback=green!3!white,
    colframe=green!40!black,
    title=Practical Example
}
\newtcolorbox{hinweisbox}{
    colback=yellow!8!white,
    colframe=yellow!50!black,
    title=Note
}
\newtcolorbox{uebungbox}{
    colback=gray!5!white,
    colframe=gray!60!black,
    title=Exercise
}

\newacronym{mbse}{MBSE}{Model-Based Systems Engineering}
\newacronym{sysml}{SysML}{Systems Modeling Language}
\newacronym{ros2}{ROS2}{Robot Operating System 2}
\newacronym{bdd}{BDD}{Block Definition Diagram (SysML v1)}
\newacronym{ibd}{IBD}{Internal Block Diagram (SysML v1)}
\newacronym{dds}{DDS}{Data Distribution Service}
\newacronym{slam}{SLAM}{Simultaneous Localization and Mapping}
\newacronym{ml}{ML}{Machine Learning}
\newacronym{ki}{AI}{Artificial Intelligence}
\newacronym{onnx}{ONNX}{Open Neural Network Exchange}
\newacronym{api}{API}{Application Programming Interface}
\newacronym{qos}{QoS}{Quality of Service}
\newacronym{req}{REQ}{Requirement}
\newacronym{tc}{TC}{Test Case}
\newacronym{fmea}{FMEA}{Failure Mode and Effects Analysis}
\newacronym{gpu}{GPU}{Graphics Processing Unit}

\newglossaryentry{requirement}{
    name=requirement,
    description={Description of a capability, property, or constraint that a system must meet}
}
\newglossaryentry{stakeholder}{
    name=stakeholder,
    description={Person, group, or organization that has an interest in the system, places requirements on it, or is affected by its effects}
}
\newglossaryentry{system}{
    name=system,
    description={A set of elements that work together for a shared purpose, within a system boundary, and with relationships to their environment}
}
\newglossaryentry{systemgrenze}{
    name=system boundary,
    description={Definition of which elements belong to the system under consideration, and which elements belong to the environment or to external systems}
}
\newglossaryentry{systemkontext}{
    name=system context,
    description={Environment of a system, including external actors, neighboring systems, constraints, and key interactions}
}
\newglossaryentry{funktion}{
    name=function,
    description={Capability or task that a system must provide, independent of which concrete component later realizes it}
}
\newglossaryentry{komponente}{
    name=component,
    description={Distinguishable part of a system, for example a sensor, a software module, a subsystem, or a mechanical part}
}
\newglossaryentry{architektur}{
    name=architecture,
    description={Underlying structure of a system: its key elements, responsibilities, interfaces, and relationships}
}
\newglossaryentry{schnittstelle}{
    name=interface,
    description={Defined transition point between parts of a system, or between the system and its environment, across which data, energy, material, or commands pass}
}
\newglossaryentry{datenfluss}{
    name=data flow,
    description={Directed exchange of information between parts of a system, for example sensor data from a camera to a processing node}
}
\newglossaryentry{usecase}{
    name=use case,
    description={Description of a goal or sequence of steps from the viewpoint of a user or external actor, for example \enquote{reach the target position}}
}
\newglossaryentry{block}{
    name=block,
    description={Structural element in SysML v1. In SysML v2, replaced by \texttt{part def} (type definition) and \texttt{part} (instance)}
}
\newglossaryentry{partdef}{
    name=part def,
    description={Type definition in SysML v2. Corresponds to the block in SysML v1. Defines which properties and parts a system element can have}
}
\newglossaryentry{part}{
    name=part,
    description={Named instance of a type defined with \texttt{part def}, inside a parent system element}
}
\newglossaryentry{port}{
    name=port,
    description={Interface point of a SysML element, used to model flows or communication relationships}
}
\newglossaryentry{connection}{
    name=connection,
    description={Link between ports or system elements that shows an exchange or a coupling between them}
}
\newglossaryentry{action}{
    name=action,
    description={In SysML v2, a behavior element within a sequence of steps; not to be confused with a ROS2 action, which is a communication pattern}
}
\newglossaryentry{state}{
    name=state,
    description={Condition of a system or part of a system in which specific rules, events, or transitions apply}
}
\newglossaryentry{transition}{
    name=transition,
    description={Change from one state to another, usually triggered by an event or a condition}
}
\newglossaryentry{guard}{
    name=guard,
    description={Condition on a branch or transition that must be true for a given path to be taken}
}
\newglossaryentry{constraint}{
    name=constraint,
    description={Restriction or rule that must hold in the model, for example a maximum braking distance or an allowed latency}
}
\newglossaryentry{calcdef}{
    name=calc def,
    description={SysML v2 definition for a calculation that structures input values, calculation logic, and result values}
}
\newglossaryentry{node}{
    name=ROS2 node,
    description={Independent process in ROS2 that handles one well-defined task, processes data, and communicates with other nodes through topics, services, or actions}
}
\newglossaryentry{topic}{
    name=topic,
    description={Message channel in ROS2 used for publish-subscribe communication}
}
\newglossaryentry{message}{
    name=message,
    description={Structure of a message in ROS2 that defines which data a topic, service, or action transmits}
}
\newglossaryentry{publisher}{
    name=publisher,
    description={ROS2 role of a node that publishes messages on a topic}
}
\newglossaryentry{subscriber}{
    name=subscriber,
    description={ROS2 role of a node that receives messages from a topic}
}
\newglossaryentry{service}{
    name=service,
    description={ROS2 communication pattern for short request-response exchanges}
}
\newglossaryentry{rosaction}{
    name=ROS2 action,
    description={ROS2 communication pattern for longer-running tasks, with a goal, feedback, a result, and an option to cancel}
}
\newglossaryentry{parameter}{
    name=parameter,
    description={Configuration value of a node or model, for example a battery threshold, a maximum speed, or a topic name}
}
\newglossaryentry{launchdatei}{
    name=launch file,
    description={File that starts several ROS2 nodes together, with their parameters and startup options}
}
\newglossaryentry{workspace}{
    name=workspace,
    description={Working directory for a ROS2 project, where source code, packages, build results, and install artifacts are organized}
}
\newglossaryentry{package}{
    name=package,
    description={Organized unit in ROS2 that groups code, configuration, interface definitions, and metadata}
}
\newglossaryentry{lifecycle-node}{
    name=lifecycle node,
    description={ROS2 node with defined states such as \texttt{unconfigured}, \texttt{inactive}, \texttt{active}, and \texttt{finalized}}
}
\newglossaryentry{traceability}{
    name=traceability,
    description={Traceable connection between requirements, architecture, implementation, and tests in the system model}
}
\newglossaryentry{verifikation}{
    name=verification,
    description={Check of whether a system or part of a system correctly meets its specified requirements}
}
\newglossaryentry{validierung}{
    name=validation,
    description={Check of whether the system that was built solves the stakeholders' actual problem and fits its intended use}
}
\newglossaryentry{testfall}{
    name=test case,
    description={Described check with a goal, a procedure, and an expected result, used to demonstrate that a requirement or function is met}
}
\newglossaryentry{impactanalysis}{
    name=impact analysis,
    description={Analysis of which requirements, components, interfaces, tests, or documents a change affects}
}
\newglossaryentry{mapping}{
    name=mapping,
    description={Assignment between elements at different levels, for example from SysML elements to ROS2 nodes, topics, and tests}
}
\newglossaryentry{middleware}{
    name=middleware,
    description={Software layer that handles communication between distributed processes (e.g., ROS2 nodes) so that they do not need to know the transport details of the other side}
}
\newglossaryentry{cyberphysical}{
    name=cyber-physical system,
    plural=cyber-physical systems,
    description={System in which physical components are controlled by software and are tightly connected to that software through sensors, actuators, or networks}
}
\newglossaryentry{conjport}{
    name=conjugated port,
    description={Port whose flow direction is reversed relative to the original port definition (marked with a tilde \textasciitilde); a sending port becomes a receiving port}
}
\newglossaryentry{sensor}{
    name=sensor,
    description={Technical element that measures a physical quantity or environment property and provides it as a signal or data value}
}
\newglossaryentry{aktor}{
    name=actuator,
    description={Technical element that turns commands into a physical effect, for example a motor, a brake mechanism, or a gripper}
}
\newglossaryentry{lidar}{
    name=lidar,
    description={Sensor principle that measures distance with laser light; commonly used in robotics for obstacle detection and mapping}
}
\newglossaryentry{sensorfusion}{
    name=sensor fusion,
    description={Combination of data from several sensors to produce a more robust or more complete picture of the environment}
}
\newglossaryentry{inferenz}{
    name=inference,
    description={Use of an already-trained machine learning model on new data, for example on a current camera image}
}
\newglossaryentry{latenz}{
    name=latency,
    description={Time delay between an event, a measurement, or an input, and the resulting system response}
}
\newglossaryentry{simulation}{
    name=simulation,
    description={Reproduction of a system or its environment, used to study behavior, algorithms, or tests without complete real hardware}
}
\newglossaryentry{modell}{
    name=model,
    description={Simplified, purpose-built description of a system that makes relevant information, relationships, and decisions visible}
}
\newglossaryentry{sicht}{
    name=view,
    description={Part of a model that shows specific information for one role or one question}
}


\title{SysML v2, MBSE, and ROS2 for Beginners\\
\large A Hands-On Textbook Based on an Autonomous Indoor Robot}
\author{Wolfgang Becker}
\date{\today}
.
% KDP's check appears to flag pages with clickable hyperlink annotations
% (cross-references, glossary terms, the whole linked table of contents) as
% suspicious regardless of their color - black link color alone did not
% remove the warnings. For print editions, hyperref is therefore loaded in
% draft mode: all \href/\ref/\gls commands remain usable and still set their
% text normally, but no longer create a PDF link annotation. The working
% edition (main_en.tex) keeps full, clickable links for on-screen use.
\makeatletter
\@ifundefined{ifdruckausgabe}{\newif\ifdruckausgabe\druckausgabefalse}{}
\makeatother
\ifdruckausgabe
    \usepackage[draft]{hyperref}
\else
    \usepackage{hyperref}
\fi

\usepackage{listings}
\usepackage{etoolbox}
\usepackage{needspace}
\usepackage{enumitem}
\usepackage{csquotes}
\usepackage{caption}
\usepackage{float}
\usepackage{tikz}
\usetikzlibrary{positioning,arrows.meta,shadows,calc,quotes,backgrounds}
\usepackage{tcolorbox}
\usepackage{amsmath}
\usepackage{makeidx}
\usepackage[acronym,toc]{glossaries}

\makeindex
\makeglossaries

\hypersetup{
    colorlinks=true,
    linkcolor=blue!50!black,
    urlcolor=blue!50!black,
    citecolor=blue!50!black,
    pdftitle={SysML v2, MBSE, and ROS2 for Beginners},
    pdfauthor={Wolfgang Becker}
}

\RedeclareSectionCommand[tocnumwidth=3.8em]{section}
\RedeclareSectionCommand[tocnumwidth=4.8em]{subsection}
\RedeclareSectionCommand[tocnumwidth=5.8em]{subsubsection}

% Section headings must not stay alone at the bottom of a page. The
% reserved space must fit the heading plus a bit of following text, or the
% start of a short code block.
\pretocmd{\section}{\Needspace{12\baselineskip}}{}{}
\pretocmd{\subsection}{\Needspace{8\baselineskip}}{}{}
\pretocmd{\subsubsection}{\Needspace{7\baselineskip}}{}{}

% Table of contents only down to section level. Subsection entries alone
% (3-digit numbers such as 1.10.1 plus a deeply indented title, not yet
% subsubsection - counting depth: chapter=0, section=1, subsection=2) already
% exceed the TOC line width on the narrower 6x9in trim and run past the
% right margin. tocdepth=2 (the previous setting) still allowed subsections,
% which was not enough - hence tocdepth=1 now. Numbering in the body text is
% unchanged (only \tocdepth, not \secnumdepth); subsections are still
% numbered, they just no longer appear in the table of contents.
\setcounter{tocdepth}{1}

\lstdefinestyle{code}{
    basicstyle=\ttfamily\small,
    breaklines=true,
    frame=single,
    numbers=left,
    numberstyle=\tiny,
    xleftmargin=0.5cm,
    columns=fullflexible,
    showstringspaces=false,
    tabsize=2,
    extendedchars=true,
    inputencoding=utf8,
    upquote=true,
    literate=
        {ä}{{\"a}}1 {ö}{{\"o}}1 {ü}{{\"u}}1
        {Ä}{{\"A}}1 {Ö}{{\"O}}1 {Ü}{{\"U}}1
        {ß}{{\ss{}}}1
}
\lstset{style=code}

% Code blocks must not start with only a few lines at the bottom of a page.
% With less than about eight code lines of space left, the whole block
% moves to the next page; longer listings may still break.
\BeforeBeginEnvironment{lstlisting}{\Needspace{10\baselineskip}}

% Style for plain text/tree displays (directory trees, flow sketches) that
% should look like \verbatim but - unlike \verbatim - must wrap long lines
% so they do not run past the text block on narrow trims (e.g. 6x9in).
\lstdefinestyle{plaintext}{
    style=code,
    numbers=none,
    frame=single,
}

\newtcolorbox{lernziele}{
    colback=blue!3!white,
    colframe=blue!50!black,
    title=Learning Objectives
}
\newtcolorbox{praxisbox}{
    colback=green!3!white,
    colframe=green!40!black,
    title=Practical Example
}
\newtcolorbox{hinweisbox}{
    colback=yellow!8!white,
    colframe=yellow!50!black,
    title=Note
}
\newtcolorbox{uebungbox}{
    colback=gray!5!white,
    colframe=gray!60!black,
    title=Exercise
}

\newacronym{mbse}{MBSE}{Model-Based Systems Engineering}
\newacronym{sysml}{SysML}{Systems Modeling Language}
\newacronym{ros2}{ROS2}{Robot Operating System 2}
\newacronym{bdd}{BDD}{Block Definition Diagram (SysML v1)}
\newacronym{ibd}{IBD}{Internal Block Diagram (SysML v1)}
\newacronym{dds}{DDS}{Data Distribution Service}
\newacronym{slam}{SLAM}{Simultaneous Localization and Mapping}
\newacronym{ml}{ML}{Machine Learning}
\newacronym{ki}{AI}{Artificial Intelligence}
\newacronym{onnx}{ONNX}{Open Neural Network Exchange}
\newacronym{api}{API}{Application Programming Interface}
\newacronym{qos}{QoS}{Quality of Service}
\newacronym{req}{REQ}{Requirement}
\newacronym{tc}{TC}{Test Case}
\newacronym{fmea}{FMEA}{Failure Mode and Effects Analysis}
\newacronym{gpu}{GPU}{Graphics Processing Unit}

\newglossaryentry{requirement}{
    name=requirement,
    description={Description of a capability, property, or constraint that a system must meet}
}
\newglossaryentry{stakeholder}{
    name=stakeholder,
    description={Person, group, or organization that has an interest in the system, places requirements on it, or is affected by its effects}
}
\newglossaryentry{system}{
    name=system,
    description={A set of elements that work together for a shared purpose, within a system boundary, and with relationships to their environment}
}
\newglossaryentry{systemgrenze}{
    name=system boundary,
    description={Definition of which elements belong to the system under consideration, and which elements belong to the environment or to external systems}
}
\newglossaryentry{systemkontext}{
    name=system context,
    description={Environment of a system, including external actors, neighboring systems, constraints, and key interactions}
}
\newglossaryentry{funktion}{
    name=function,
    description={Capability or task that a system must provide, independent of which concrete component later realizes it}
}
\newglossaryentry{komponente}{
    name=component,
    description={Distinguishable part of a system, for example a sensor, a software module, a subsystem, or a mechanical part}
}
\newglossaryentry{architektur}{
    name=architecture,
    description={Underlying structure of a system: its key elements, responsibilities, interfaces, and relationships}
}
\newglossaryentry{schnittstelle}{
    name=interface,
    description={Defined transition point between parts of a system, or between the system and its environment, across which data, energy, material, or commands pass}
}
\newglossaryentry{datenfluss}{
    name=data flow,
    description={Directed exchange of information between parts of a system, for example sensor data from a camera to a processing node}
}
\newglossaryentry{usecase}{
    name=use case,
    description={Description of a goal or sequence of steps from the viewpoint of a user or external actor, for example \enquote{reach the target position}}
}
\newglossaryentry{block}{
    name=block,
    description={Structural element in SysML v1. In SysML v2, replaced by \texttt{part def} (type definition) and \texttt{part} (instance)}
}
\newglossaryentry{partdef}{
    name=part def,
    description={Type definition in SysML v2. Corresponds to the block in SysML v1. Defines which properties and parts a system element can have}
}
\newglossaryentry{part}{
    name=part,
    description={Named instance of a type defined with \texttt{part def}, inside a parent system element}
}
\newglossaryentry{port}{
    name=port,
    description={Interface point of a SysML element, used to model flows or communication relationships}
}
\newglossaryentry{connection}{
    name=connection,
    description={Link between ports or system elements that shows an exchange or a coupling between them}
}
\newglossaryentry{action}{
    name=action,
    description={In SysML v2, a behavior element within a sequence of steps; not to be confused with a ROS2 action, which is a communication pattern}
}
\newglossaryentry{state}{
    name=state,
    description={Condition of a system or part of a system in which specific rules, events, or transitions apply}
}
\newglossaryentry{transition}{
    name=transition,
    description={Change from one state to another, usually triggered by an event or a condition}
}
\newglossaryentry{guard}{
    name=guard,
    description={Condition on a branch or transition that must be true for a given path to be taken}
}
\newglossaryentry{constraint}{
    name=constraint,
    description={Restriction or rule that must hold in the model, for example a maximum braking distance or an allowed latency}
}
\newglossaryentry{calcdef}{
    name=calc def,
    description={SysML v2 definition for a calculation that structures input values, calculation logic, and result values}
}
\newglossaryentry{node}{
    name=ROS2 node,
    description={Independent process in ROS2 that handles one well-defined task, processes data, and communicates with other nodes through topics, services, or actions}
}
\newglossaryentry{topic}{
    name=topic,
    description={Message channel in ROS2 used for publish-subscribe communication}
}
\newglossaryentry{message}{
    name=message,
    description={Structure of a message in ROS2 that defines which data a topic, service, or action transmits}
}
\newglossaryentry{publisher}{
    name=publisher,
    description={ROS2 role of a node that publishes messages on a topic}
}
\newglossaryentry{subscriber}{
    name=subscriber,
    description={ROS2 role of a node that receives messages from a topic}
}
\newglossaryentry{service}{
    name=service,
    description={ROS2 communication pattern for short request-response exchanges}
}
\newglossaryentry{rosaction}{
    name=ROS2 action,
    description={ROS2 communication pattern for longer-running tasks, with a goal, feedback, a result, and an option to cancel}
}
\newglossaryentry{parameter}{
    name=parameter,
    description={Configuration value of a node or model, for example a battery threshold, a maximum speed, or a topic name}
}
\newglossaryentry{launchdatei}{
    name=launch file,
    description={File that starts several ROS2 nodes together, with their parameters and startup options}
}
\newglossaryentry{workspace}{
    name=workspace,
    description={Working directory for a ROS2 project, where source code, packages, build results, and install artifacts are organized}
}
\newglossaryentry{package}{
    name=package,
    description={Organized unit in ROS2 that groups code, configuration, interface definitions, and metadata}
}
\newglossaryentry{lifecycle-node}{
    name=lifecycle node,
    description={ROS2 node with defined states such as \texttt{unconfigured}, \texttt{inactive}, \texttt{active}, and \texttt{finalized}}
}
\newglossaryentry{traceability}{
    name=traceability,
    description={Traceable connection between requirements, architecture, implementation, and tests in the system model}
}
\newglossaryentry{verifikation}{
    name=verification,
    description={Check of whether a system or part of a system correctly meets its specified requirements}
}
\newglossaryentry{validierung}{
    name=validation,
    description={Check of whether the system that was built solves the stakeholders' actual problem and fits its intended use}
}
\newglossaryentry{testfall}{
    name=test case,
    description={Described check with a goal, a procedure, and an expected result, used to demonstrate that a requirement or function is met}
}
\newglossaryentry{impactanalysis}{
    name=impact analysis,
    description={Analysis of which requirements, components, interfaces, tests, or documents a change affects}
}
\newglossaryentry{mapping}{
    name=mapping,
    description={Assignment between elements at different levels, for example from SysML elements to ROS2 nodes, topics, and tests}
}
\newglossaryentry{middleware}{
    name=middleware,
    description={Software layer that handles communication between distributed processes (e.g., ROS2 nodes) so that they do not need to know the transport details of the other side}
}
\newglossaryentry{cyberphysical}{
    name=cyber-physical system,
    plural=cyber-physical systems,
    description={System in which physical components are controlled by software and are tightly connected to that software through sensors, actuators, or networks}
}
\newglossaryentry{conjport}{
    name=conjugated port,
    description={Port whose flow direction is reversed relative to the original port definition (marked with a tilde \textasciitilde); a sending port becomes a receiving port}
}
\newglossaryentry{sensor}{
    name=sensor,
    description={Technical element that measures a physical quantity or environment property and provides it as a signal or data value}
}
\newglossaryentry{aktor}{
    name=actuator,
    description={Technical element that turns commands into a physical effect, for example a motor, a brake mechanism, or a gripper}
}
\newglossaryentry{lidar}{
    name=lidar,
    description={Sensor principle that measures distance with laser light; commonly used in robotics for obstacle detection and mapping}
}
\newglossaryentry{sensorfusion}{
    name=sensor fusion,
    description={Combination of data from several sensors to produce a more robust or more complete picture of the environment}
}
\newglossaryentry{inferenz}{
    name=inference,
    description={Use of an already-trained machine learning model on new data, for example on a current camera image}
}
\newglossaryentry{latenz}{
    name=latency,
    description={Time delay between an event, a measurement, or an input, and the resulting system response}
}
\newglossaryentry{simulation}{
    name=simulation,
    description={Reproduction of a system or its environment, used to study behavior, algorithms, or tests without complete real hardware}
}
\newglossaryentry{modell}{
    name=model,
    description={Simplified, purpose-built description of a system that makes relevant information, relationships, and decisions visible}
}
\newglossaryentry{sicht}{
    name=view,
    description={Part of a model that shows specific information for one role or one question}
}


\title{SysML v2, MBSE, and ROS2 for Beginners\\
\large A Hands-On Textbook Based on an Autonomous Indoor Robot}
\author{Wolfgang Becker}
\date{\today}
.
% KDP's check appears to flag pages with clickable hyperlink annotations
% (cross-references, glossary terms, the whole linked table of contents) as
% suspicious regardless of their color - black link color alone did not
% remove the warnings. For print editions, hyperref is therefore loaded in
% draft mode: all \href/\ref/\gls commands remain usable and still set their
% text normally, but no longer create a PDF link annotation. The working
% edition (main_en.tex) keeps full, clickable links for on-screen use.
\makeatletter
\@ifundefined{ifdruckausgabe}{\newif\ifdruckausgabe\druckausgabefalse}{}
\makeatother
\ifdruckausgabe
    \usepackage[draft]{hyperref}
\else
    \usepackage{hyperref}
\fi

\usepackage{listings}
\usepackage{etoolbox}
\usepackage{needspace}
\usepackage{enumitem}
\usepackage{csquotes}
\usepackage{caption}
\usepackage{float}
\usepackage{tikz}
\usetikzlibrary{positioning,arrows.meta,shadows,calc,quotes,backgrounds}
\usepackage{tcolorbox}
\usepackage{amsmath}
\usepackage{makeidx}
\usepackage[acronym,toc]{glossaries}

\makeindex
\makeglossaries

\hypersetup{
    colorlinks=true,
    linkcolor=blue!50!black,
    urlcolor=blue!50!black,
    citecolor=blue!50!black,
    pdftitle={SysML v2, MBSE, and ROS2 for Beginners},
    pdfauthor={Wolfgang Becker}
}

\RedeclareSectionCommand[tocnumwidth=3.8em]{section}
\RedeclareSectionCommand[tocnumwidth=4.8em]{subsection}
\RedeclareSectionCommand[tocnumwidth=5.8em]{subsubsection}

% Section headings must not stay alone at the bottom of a page. The
% reserved space must fit the heading plus a bit of following text, or the
% start of a short code block.
\pretocmd{\section}{\Needspace{12\baselineskip}}{}{}
\pretocmd{\subsection}{\Needspace{8\baselineskip}}{}{}
\pretocmd{\subsubsection}{\Needspace{7\baselineskip}}{}{}

% Table of contents only down to section level. Subsection entries alone
% (3-digit numbers such as 1.10.1 plus a deeply indented title, not yet
% subsubsection - counting depth: chapter=0, section=1, subsection=2) already
% exceed the TOC line width on the narrower 6x9in trim and run past the
% right margin. tocdepth=2 (the previous setting) still allowed subsections,
% which was not enough - hence tocdepth=1 now. Numbering in the body text is
% unchanged (only \tocdepth, not \secnumdepth); subsections are still
% numbered, they just no longer appear in the table of contents.
\setcounter{tocdepth}{1}

\lstdefinestyle{code}{
    basicstyle=\ttfamily\small,
    breaklines=true,
    frame=single,
    numbers=left,
    numberstyle=\tiny,
    xleftmargin=0.5cm,
    columns=fullflexible,
    showstringspaces=false,
    tabsize=2,
    extendedchars=true,
    inputencoding=utf8,
    upquote=true,
    literate=
        {ä}{{\"a}}1 {ö}{{\"o}}1 {ü}{{\"u}}1
        {Ä}{{\"A}}1 {Ö}{{\"O}}1 {Ü}{{\"U}}1
        {ß}{{\ss{}}}1
}
\lstset{style=code}

% Code blocks must not start with only a few lines at the bottom of a page.
% With less than about eight code lines of space left, the whole block
% moves to the next page; longer listings may still break.
\BeforeBeginEnvironment{lstlisting}{\Needspace{10\baselineskip}}

% Style for plain text/tree displays (directory trees, flow sketches) that
% should look like \verbatim but - unlike \verbatim - must wrap long lines
% so they do not run past the text block on narrow trims (e.g. 6x9in).
\lstdefinestyle{plaintext}{
    style=code,
    numbers=none,
    frame=single,
}

\newtcolorbox{lernziele}{
    colback=blue!3!white,
    colframe=blue!50!black,
    title=Learning Objectives
}
\newtcolorbox{praxisbox}{
    colback=green!3!white,
    colframe=green!40!black,
    title=Practical Example
}
\newtcolorbox{hinweisbox}{
    colback=yellow!8!white,
    colframe=yellow!50!black,
    title=Note
}
\newtcolorbox{uebungbox}{
    colback=gray!5!white,
    colframe=gray!60!black,
    title=Exercise
}

\newacronym{mbse}{MBSE}{Model-Based Systems Engineering}
\newacronym{sysml}{SysML}{Systems Modeling Language}
\newacronym{ros2}{ROS2}{Robot Operating System 2}
\newacronym{bdd}{BDD}{Block Definition Diagram (SysML v1)}
\newacronym{ibd}{IBD}{Internal Block Diagram (SysML v1)}
\newacronym{dds}{DDS}{Data Distribution Service}
\newacronym{slam}{SLAM}{Simultaneous Localization and Mapping}
\newacronym{ml}{ML}{Machine Learning}
\newacronym{ki}{AI}{Artificial Intelligence}
\newacronym{onnx}{ONNX}{Open Neural Network Exchange}
\newacronym{api}{API}{Application Programming Interface}
\newacronym{qos}{QoS}{Quality of Service}
\newacronym{req}{REQ}{Requirement}
\newacronym{tc}{TC}{Test Case}
\newacronym{fmea}{FMEA}{Failure Mode and Effects Analysis}
\newacronym{gpu}{GPU}{Graphics Processing Unit}

\newglossaryentry{requirement}{
    name=requirement,
    description={Description of a capability, property, or constraint that a system must meet}
}
\newglossaryentry{stakeholder}{
    name=stakeholder,
    description={Person, group, or organization that has an interest in the system, places requirements on it, or is affected by its effects}
}
\newglossaryentry{system}{
    name=system,
    description={A set of elements that work together for a shared purpose, within a system boundary, and with relationships to their environment}
}
\newglossaryentry{systemgrenze}{
    name=system boundary,
    description={Definition of which elements belong to the system under consideration, and which elements belong to the environment or to external systems}
}
\newglossaryentry{systemkontext}{
    name=system context,
    description={Environment of a system, including external actors, neighboring systems, constraints, and key interactions}
}
\newglossaryentry{funktion}{
    name=function,
    description={Capability or task that a system must provide, independent of which concrete component later realizes it}
}
\newglossaryentry{komponente}{
    name=component,
    description={Distinguishable part of a system, for example a sensor, a software module, a subsystem, or a mechanical part}
}
\newglossaryentry{architektur}{
    name=architecture,
    description={Underlying structure of a system: its key elements, responsibilities, interfaces, and relationships}
}
\newglossaryentry{schnittstelle}{
    name=interface,
    description={Defined transition point between parts of a system, or between the system and its environment, across which data, energy, material, or commands pass}
}
\newglossaryentry{datenfluss}{
    name=data flow,
    description={Directed exchange of information between parts of a system, for example sensor data from a camera to a processing node}
}
\newglossaryentry{usecase}{
    name=use case,
    description={Description of a goal or sequence of steps from the viewpoint of a user or external actor, for example \enquote{reach the target position}}
}
\newglossaryentry{block}{
    name=block,
    description={Structural element in SysML v1. In SysML v2, replaced by \texttt{part def} (type definition) and \texttt{part} (instance)}
}
\newglossaryentry{partdef}{
    name=part def,
    description={Type definition in SysML v2. Corresponds to the block in SysML v1. Defines which properties and parts a system element can have}
}
\newglossaryentry{part}{
    name=part,
    description={Named instance of a type defined with \texttt{part def}, inside a parent system element}
}
\newglossaryentry{port}{
    name=port,
    description={Interface point of a SysML element, used to model flows or communication relationships}
}
\newglossaryentry{connection}{
    name=connection,
    description={Link between ports or system elements that shows an exchange or a coupling between them}
}
\newglossaryentry{action}{
    name=action,
    description={In SysML v2, a behavior element within a sequence of steps; not to be confused with a ROS2 action, which is a communication pattern}
}
\newglossaryentry{state}{
    name=state,
    description={Condition of a system or part of a system in which specific rules, events, or transitions apply}
}
\newglossaryentry{transition}{
    name=transition,
    description={Change from one state to another, usually triggered by an event or a condition}
}
\newglossaryentry{guard}{
    name=guard,
    description={Condition on a branch or transition that must be true for a given path to be taken}
}
\newglossaryentry{constraint}{
    name=constraint,
    description={Restriction or rule that must hold in the model, for example a maximum braking distance or an allowed latency}
}
\newglossaryentry{calcdef}{
    name=calc def,
    description={SysML v2 definition for a calculation that structures input values, calculation logic, and result values}
}
\newglossaryentry{node}{
    name=ROS2 node,
    description={Independent process in ROS2 that handles one well-defined task, processes data, and communicates with other nodes through topics, services, or actions}
}
\newglossaryentry{topic}{
    name=topic,
    description={Message channel in ROS2 used for publish-subscribe communication}
}
\newglossaryentry{message}{
    name=message,
    description={Structure of a message in ROS2 that defines which data a topic, service, or action transmits}
}
\newglossaryentry{publisher}{
    name=publisher,
    description={ROS2 role of a node that publishes messages on a topic}
}
\newglossaryentry{subscriber}{
    name=subscriber,
    description={ROS2 role of a node that receives messages from a topic}
}
\newglossaryentry{service}{
    name=service,
    description={ROS2 communication pattern for short request-response exchanges}
}
\newglossaryentry{rosaction}{
    name=ROS2 action,
    description={ROS2 communication pattern for longer-running tasks, with a goal, feedback, a result, and an option to cancel}
}
\newglossaryentry{parameter}{
    name=parameter,
    description={Configuration value of a node or model, for example a battery threshold, a maximum speed, or a topic name}
}
\newglossaryentry{launchdatei}{
    name=launch file,
    description={File that starts several ROS2 nodes together, with their parameters and startup options}
}
\newglossaryentry{workspace}{
    name=workspace,
    description={Working directory for a ROS2 project, where source code, packages, build results, and install artifacts are organized}
}
\newglossaryentry{package}{
    name=package,
    description={Organized unit in ROS2 that groups code, configuration, interface definitions, and metadata}
}
\newglossaryentry{lifecycle-node}{
    name=lifecycle node,
    description={ROS2 node with defined states such as \texttt{unconfigured}, \texttt{inactive}, \texttt{active}, and \texttt{finalized}}
}
\newglossaryentry{traceability}{
    name=traceability,
    description={Traceable connection between requirements, architecture, implementation, and tests in the system model}
}
\newglossaryentry{verifikation}{
    name=verification,
    description={Check of whether a system or part of a system correctly meets its specified requirements}
}
\newglossaryentry{validierung}{
    name=validation,
    description={Check of whether the system that was built solves the stakeholders' actual problem and fits its intended use}
}
\newglossaryentry{testfall}{
    name=test case,
    description={Described check with a goal, a procedure, and an expected result, used to demonstrate that a requirement or function is met}
}
\newglossaryentry{impactanalysis}{
    name=impact analysis,
    description={Analysis of which requirements, components, interfaces, tests, or documents a change affects}
}
\newglossaryentry{mapping}{
    name=mapping,
    description={Assignment between elements at different levels, for example from SysML elements to ROS2 nodes, topics, and tests}
}
\newglossaryentry{middleware}{
    name=middleware,
    description={Software layer that handles communication between distributed processes (e.g., ROS2 nodes) so that they do not need to know the transport details of the other side}
}
\newglossaryentry{cyberphysical}{
    name=cyber-physical system,
    plural=cyber-physical systems,
    description={System in which physical components are controlled by software and are tightly connected to that software through sensors, actuators, or networks}
}
\newglossaryentry{conjport}{
    name=conjugated port,
    description={Port whose flow direction is reversed relative to the original port definition (marked with a tilde \textasciitilde); a sending port becomes a receiving port}
}
\newglossaryentry{sensor}{
    name=sensor,
    description={Technical element that measures a physical quantity or environment property and provides it as a signal or data value}
}
\newglossaryentry{aktor}{
    name=actuator,
    description={Technical element that turns commands into a physical effect, for example a motor, a brake mechanism, or a gripper}
}
\newglossaryentry{lidar}{
    name=lidar,
    description={Sensor principle that measures distance with laser light; commonly used in robotics for obstacle detection and mapping}
}
\newglossaryentry{sensorfusion}{
    name=sensor fusion,
    description={Combination of data from several sensors to produce a more robust or more complete picture of the environment}
}
\newglossaryentry{inferenz}{
    name=inference,
    description={Use of an already-trained machine learning model on new data, for example on a current camera image}
}
\newglossaryentry{latenz}{
    name=latency,
    description={Time delay between an event, a measurement, or an input, and the resulting system response}
}
\newglossaryentry{simulation}{
    name=simulation,
    description={Reproduction of a system or its environment, used to study behavior, algorithms, or tests without complete real hardware}
}
\newglossaryentry{modell}{
    name=model,
    description={Simplified, purpose-built description of a system that makes relevant information, relationships, and decisions visible}
}
\newglossaryentry{sicht}{
    name=view,
    description={Part of a model that shows specific information for one role or one question}
}


\title{SysML v2, MBSE, and ROS2 for Beginners\\
\large A Hands-On Textbook Based on an Autonomous Indoor Robot}
\author{Wolfgang Becker}
\date{\today}
